\begingroup\let\clearpageforsection\relax\reviewexercisesheader{}\endgroup

% 27

\eoce{\qt{Ujian susulan\label{makeup_exam}} Dalam sebuah kelas yang terdiri atas
25 mahasiswa, 24 mahasiswa mengikuti ujian di kelas dan 1 mahasiswa mengikuti
ujian susulan pada hari berikutnya. Dosen menilai 24 ujian pertama dan
memperoleh rerata nilai 74 poin dengan simpangan baku 8.9 poin. Mahasiswa
yang mengikuti ujian susulan pada hari berikutnya memperoleh nilai 64 poin.
\begin{parts}
\item Apakah nilai mahasiswa yang mengikuti ujian susulan tersebut menaikkan
atau menurunkan rerata nilai?
\item Berapakah rerata nilai yang baru?
\item Apakah nilai mahasiswa yang mengikuti ujian susulan tersebut menaikkan
atau menurunkan simpangan baku nilai-nilai itu?
\end{parts}
}{}

% 28

\eoce{\qt{Kematian bayi\label{infant_mortality}} Angka kematian bayi
didefinisikan sebagai banyaknya kematian bayi per 1,000 kelahiran hidup. Angka
ini sering digunakan sebagai indikator tingkat kesehatan suatu negara.
Histogram frekuensi relatif di bawah ini menunjukkan distribusi perkiraan angka
kematian bayi untuk 224 negara yang datanya tersedia pada 2014.
\footfullcite{data:ciaFactbook}

\noindent\begin{minipage}[c]{0.43\textwidth}
\begin{parts}
\item Perkirakan Q1, median, dan Q3 dari histogram tersebut.
\item Apakah Anda memperkirakan rata-rata kumpulan data ini lebih kecil atau
lebih besar daripada mediannya? Jelaskan penalaran Anda.
\end{parts} \vfill \
\end{minipage}
\begin{minipage}[c]{0.52\textwidth}
\hfill%
\Figures[Sebuah histogram ditampilkan untuk variabel "Kematian Bayi (per 1000
kelahiran hidup)" dengan rentang sumbu dari 0 hingga 120. Sumbu vertikal
histogram menyatakan "Proporsi Negara" dan membentang dari 0 hingga 0.4.
Untuk setiap interval, jumlah negara, tinggi proporsi eksak, dan tinggi yang
dibulatkan ke dua tempat desimal adalah sebagai berikut: 0 hingga 10: 84 negara
(84/224; 0.38); 10 hingga 20: 50 negara (50/224; 0.22); 20 hingga 30: 24 negara
(24/224; 0.11); 30 hingga 40: 13 negara (13/224; 0.06); 40 hingga 50: 14 negara
(14/224; 0.06); 50 hingga 60: 17 negara (17/224; 0.08); 60 hingga 70: 8 negara
(8/224; 0.04); 70 hingga 80: 7 negara (7/224; 0.03); 80 hingga 90: 1 negara
(1/224; 0.00); 90 hingga 100: 3 negara (3/224; 0.01); 100 hingga 110: 2 negara
(2/224; 0.01); dan 110 hingga 120: 1 negara (1/224; 0.00).]{0.85}{eoce/infant_mortality_rel_freq}{infant_mortality_rel_freq_hist}
\end{minipage}
}{}

% 29

\eoce{\qt{Penonton televisi\label{dist_shape_TV_watchers}} Para siswa dalam
sebuah kelas AP Statistics ditanya berapa jam mereka menonton televisi setiap
minggu (termasuk melalui layanan streaming daring). Sampel ini menghasilkan
rerata 4.71 jam dengan simpangan baku 4.18 jam. Apakah distribusi banyaknya
jam menonton televisi per minggu tersebut simetris? Jika tidak, bentuk seperti
apa yang Anda perkirakan untuk distribusi ini? Jelaskan penalaran Anda.
}{}

% 30

\eoce{\qt{Statistik baru\label{new_stat}} Statistik
$\frac{\bar{x}}{median}$ dapat digunakan sebagai ukuran kemencengan. Misalkan
kita memiliki distribusi yang seluruh pengamatannya lebih besar daripada 0,
$x_i > 0$. Bagaimanakah bentuk distribusi yang diperkirakan dalam setiap kondisi
berikut? Jelaskan penalaran Anda.
\begin{parts}
\item $\frac{\bar{x}}{median} = 1$
\item $\frac{\bar{x}}{median} < 1$
\item $\frac{\bar{x}}{median} > 1$
\end{parts}
}{}

% 31

\eoce{\qt{Pemenang Oscar\label{oscar_winners}} Penghargaan Oscar untuk kategori
Aktor Terbaik dan Aktris Terbaik pertama kali diberikan pada 1929. Histogram
di bawah ini menunjukkan distribusi usia seluruh pemenang kategori Aktor
Terbaik dan Aktris Terbaik dari 1929 hingga 2018. Statistik ringkasan untuk
kedua distribusi ini juga disediakan. Bandingkan distribusi usia pemenang
kategori Aktor Terbaik dan Aktris Terbaik.\footfullcite{data:oscars} \\
\begin{minipage}[c]{0.72\textwidth}
\begin{center}
\Figures[Ditampilkan dua histogram, satu untuk "Aktris Terbaik" dan satu lagi
untuk "Aktor Terbaik", dengan nilai pada kedua histogram membentang dari 15
hingga 85. Tinggi interval-interval kelas pada histogram Aktris Terbaik adalah
sebagai berikut: interval 15 hingga 25 memiliki tinggi 9, 25 hingga 35 memiliki
tinggi 50, 35 hingga 45 memiliki tinggi 19, 45 hingga 55 memiliki tinggi 6, 55
hingga 65 memiliki tinggi 8, 65 hingga 75 memiliki tinggi 1, dan 75 hingga 85
memiliki tinggi 1. Tinggi interval-interval kelas pada histogram Aktor Terbaik
adalah sebagai berikut: interval 15 hingga 25 memiliki tinggi 0, 25 hingga 35
memiliki tinggi 14, 35 hingga 45 memiliki tinggi 45, 45 hingga 55 memiliki
tinggi 23, 55 hingga 65 memiliki tinggi 11, 65 hingga 75 memiliki tinggi 0, dan
75 hingga 85 memiliki tinggi 1.]{0.95}{eoce/oscar_winners}{oscars_winners_hist}
\end{center}
\end{minipage}
\begin{minipage}[c]{0.27\textwidth}
{\small
\begin{tabular}{l c}
\hline
             & Aktris Terbaik  \\
\hline
Rata-rata    & 36.2      \\
Simp. baku   & 11.9      \\
n            & 92        \\  
             & \\
             & \\
             & \\
             & \\
             & \\
\hline
             & Aktor Terbaik \\
\hline
Rata-rata    & 43.8 \\
Simp. baku   & 8.83 \\
n            & 92
\end{tabular}
}
\end{minipage}
}{}

% 32

\eoce{\qt{Nilai ujian\label{dist_shape_exam_scores}} Rerata nilai sebuah
ujian sejarah (dengan nilai maksimum 100 poin) adalah 85, dengan simpangan baku
15. Apakah distribusi nilai ujian ini simetris? Jika tidak, bentuk seperti apa
yang Anda perkirakan untuk distribusi tersebut? Jelaskan penalaran Anda.
}{}

% 33

\eoce{\qt{Nilai statistika\label{stats_scores_box}} Berikut adalah nilai ujian
akhir dua puluh mahasiswa yang mengikuti mata kuliah statistika pengantar.
\begin{center}
57, 66, 69, 71, 72, 73, 74, 77, 78, 78, 79, 79, 81, 81, 82, 83, 83, 88, 89, 94
\end{center}
Buatlah diagram kotak untuk distribusi nilai tersebut. Ringkasan lima angka di
bawah ini mungkin berguna.
\begin{center}
\renewcommand\arraystretch{1.5}
\begin{tabular}{ccccc}
Min & Q1    & Q2 (Median)   & Q3    & Maks \\
\hline
57  & 72.5  & 78.5          & 82.5  & 94 \\
\end{tabular}
\end{center}
}{}

% 34

\eoce{\qt{Pemenang maraton\label{marathon_winners}} Histogram dan diagram kotak
di bawah ini menunjukkan distribusi waktu finis dalam jam bagi pemenang pria
dan wanita Maraton New York antara 1970 dan 1999.
\begin{center}
\Figures[Ditampilkan dua grafik, yaitu sebuah histogram dan sebuah diagram
kotak, dengan rentang sumbu masing-masing dari 2.0 hingga 3.2. Interval-interval
kelas pada histogram adalah sebagai berikut: interval 2.0 hingga 2.2 memiliki
tinggi 21, 2.2 hingga 2.4 memiliki tinggi 6, 2.4 hingga 2.6 memiliki tinggi 25,
2.6 hingga 2.8 memiliki tinggi 3, 2.8 hingga 3.0 memiliki tinggi 2, dan 3.0
hingga 3.2 memiliki tinggi 2. Pada diagram kotak, kotak membentang dari 2.2
hingga 2.5 dengan garis median di 2.4. Kumis membentang dari sekitar 2.15 hingga
2.75. Empat titik ditandai di atas kumis atas pada 2.9, 3.0, 3.10, dan
3.15.]{0.56}{eoce/marathon_winners}{marathon_winners_hist_box}
\end{center}
\begin{parts}
\item Ciri-ciri distribusi apa yang tampak pada histogram tetapi tidak pada
diagram kotak? Ciri-ciri apa yang tampak pada diagram kotak tetapi tidak pada
histogram?
\item Apa yang mungkin menyebabkan distribusi bimodal tersebut? Jelaskan.
\newpage
\vspace*{\fill}
\item Bandingkan distribusi waktu maraton pria dan wanita berdasarkan diagram
kotak di bawah ini.
\begin{center}
\Figures[Ditampilkan diagram kotak berdampingan untuk waktu lari maraton, satu
diagram kotak untuk pria dan satu untuk wanita. Sumbu waktu lari membentang dari
2.0 hingga 3.2. Semua nilai yang diuraikan berikut merupakan perkiraan. Pada
diagram kotak pria, kotak membentang dari 2.16 hingga 2.22 dengan garis median
pada 2.19. Kumis membentang dari 2.12 hingga 2.27. Terdapat 6 titik di atas
kumis atas pada 2.32, 2.36, 2.38, 2.44, 2.46, dan 2.50. Pada diagram kotak wanita,
kotak membentang dari 2.44 hingga 2.52 dengan median 2.46. Kumis membentang dari
2.41 hingga 2.57. Terdapat 6 titik di atas kumis atas pada 2.72, 2.78, 2.9,
2.92, 3.12, dan 3.15.]{0.90}{eoce/marathon_winners}{marathon_winners_gender_box}
\end{center}
\item Grafik deret waktu di bawah ini merupakan cara lain untuk menelaah data
tersebut. Uraikan apa yang tampak pada grafik ini tetapi tidak pada grafik lain.
\end{parts}
\begin{center}
\Figures[Ditampilkan sebuah grafik deret waktu yang dalam kasus ini tampak
seperti diagram pencar. Sumbu horizontal menyatakan tahun dan membentang dari
1970 hingga 2000, sedangkan sumbu vertikal menyatakan "Waktu maraton" dan
membentang dari 2.0 hingga 3.2 jam. Penanda wanita berupa tanda silang dan
penanda pria berupa titik; warnanya juga berbeda. Terdapat satu titik pria
untuk setiap tahun dan, kecuali pada tahun pertama, satu titik wanita. Seri
pria dimulai pada sekitar 2.5 pada 1970, sedangkan seri wanita dimulai pada
sekitar 2.9 pada 1971. Waktu finis kedua kelompok naik turun selama beberapa
tahun, lalu menurun pada 1975 atau 1976 menjadi 2.2 untuk pria dan 2.7 untuk
wanita. Waktu finis pemenang wanita terus menurun selama beberapa tahun hingga
sekitar 2.5. Sepanjang tahun-tahun berikutnya, waktu finis berfluktuasi naik
atau turun sekitar 0.1 jam dari tahun ke tahun, tetapi relatif stabil hingga
1999, tahun terakhir yang datanya ditampilkan.]{0.95}{eoce/marathon_winners}{marathon_winners_time_series} \\
\end{center}
\vspace*{\fill}
}{}
